\section{n1. Independent whole416 canonical/surface reconstruction; exact smooth source theorem and source qualifications preserved.}
\textit{Supports: G331.}

Relative to the ENTIRE qualified G331 only, this note reconstructs the canonical envelope and the dimension-two normal result used below. Current416 and current413 notes, source ledgers, programs, verdicts and compiler records are DATA, not accepted supports. Every G331 annotation remains binding: its core consists of actual negative effective finite REAL Cartier presentations on directed ACTUAL primary blowups, with nonnegative degree on EVERY contracted proper integral curve; its C and D are the uniform and actual decreasing-net completions. Its comparison with ideal classes and full-fibre positive descent are granted. Arbitrary birational-model cofinality, envelopes and characteristic-free resolution are not granted.

Let R be excellent reduced local, a proper primary gauge, and X=V\_a(R), retaining EVERY prime support. For finite continuous signed f put A\_0=R and A\_m=\{g in R: v(g)+m f(v)\textgreater{}=0 for every v in X\}. If M=max(0,-min f), then a\textasciicircum{}N is contained in A\_m for N\textgreater{}=mM, N\textgreater{}0. Addition, scalar multiplication and the ultrametric inequality prove idealness. A negative value of f excludes every unit, making A\_m proper primary; if f\textgreater{}=0, A\_m=R. Multiplication gives A\_m A\_n subset A\_(m+n). A monic ideal-integrality equation proves A\_m integrally closed by the unique least-valued leading-term argument, including infinite values by treating them separately. The profile h\_m=-v(A\_m)/m is continuous and below f. Its divisor -E\_(A\_m)/m on Bl\_(A\_m)(R) is ACTUAL negative effective Cartier and relatively ample. Its restriction to every contracted integral curve is ample, hence has positive degree. The unit case is zero. Thus the approximants have independently geometric positivity before their limit is identified.

The raw pointwise supremum p(f)=sup\{u in D:u\textless{}=f\} equals the supremum over C, the supremum of rational primary/unit competitors, and sup\_m h\_m. Indeed a rational competitor -v(I)/r has I subset A\_r. Uniform rational approximation of any C competitor, followed by a small negative rational gauge constant (ideal multiplication after common denominators), yields rational competitors converging from below. If u=inf\_i g\_i\textless{}=f is an actual decreasing C-net, the open sets \{g\_i\textless{}f+epsilon\} cover compact X and increase under directed refinement. A finite subcover and one common index give g\_i-epsilon\textless{}=f everywhere and u\textless{}=p\_C(f)+epsilon. This also covers bottom and partial poles. No positivity of a raw supremum is assumed.

Put M\_1=||v(A\_1)|| and H\_N=max\_(1\textless{}=r\textless{}=N) h\_r. For m=kr+s, 0\textless{}=s\textless{}r, grading gives h\_m\textgreater{}=(kr/m)h\_r-sM\_1/m\textgreater{}=h\_r-rM\_1/m. Therefore H\_N-NM\_1/m\textless{}=h\_m\textless{}=p. All degrees converge pointwise to a finite nonpositive lower semicontinuous p. Continuity of p is equivalent to uniform convergence of H\_N, and to uniform convergence of h\_m: if p is continuous, the increasing open sets \{H\_N\textgreater{}p-epsilon\} cover compact X; the displayed bound then controls every sufficiently large degree. Conversely a uniform limit of h\_m is continuous and belongs to closed geometric C. In this case p is the greatest C/D minorant. Subtracting an epsilon gauge constant from a D competitor proves ||p(f)-p(g)||\textless{}=||f-g||. A finite homogeneous generating set for the canonical algebra would make p the rational profile -v((g\_i\textasciicircum{}(L/d\_i)))/L, for L divisible by all generator degrees d\_i; evaluate homogeneous expressions and pure powers. This elementary forward implication needs no section-algebra finite-generation theorem.

Primary change b is exact: d(v)=v(b)\textgreater{}0, T(v)=v/d(v), and f\_b(Tv)=f\_a(v)/d(v). Every defining inequality divides by d, so ALL A\_m are unchanged and profiles/envelopes divide by d. Both d and 1/d are bounded. G331's completion homeomorphism and exact primary contraction transport the canonical ideals and profiles. If f\_i decrease through continuous obstacles with C envelopes, inf\_i p(f\_i) is an actual decreasing C-net and is the greatest D minorant of inf\_i f\_i, since every competitor is below each p(f\_i).

Here is the repaired dimension-two normal proof, without importing the current416 verdict. For an excellent normal local DOMAIN R of dimension TWO, dominate any finite primary-ideal list INCLUDING a by its product blowup and finite normalization. Its punctured base is regular, and the normal projective surface has finitely many singular points: excellence gives an open regular locus and normality excludes codimension-one singularities. Their local completions are normal by excellence. These verify EXACTLY condition (4) in Stacks54.14.5. The complete immediate proof, with54.14.3--4, resolves this surface by finitely many normalized blowups over its singular closed points. Deeper complete-normal-surface resolution, rational/double-point reduction, completion transfer and gluing remain the disclosed published foundations; they are not a fresh recursive audit or a higher-dimensional resolution premise.

The resulting regular Y is projective and unchanged off the marked point. Pullback of the initial negative tautological bundle is ample under finite normalization; the successive negative tautological bundles and sufficiently large pullback powers combine by Stacks29.38.7--8 into O(-E), where E is effective exceptional integral with full support. The initial primary divisor ensures positivity at EVERY exceptional prime. Stacks37.53.6 gives H\textasciicircum{}0(Y,O)=R: Y is integral proper birational over normal R and its generic fibre has sections Frac(R).

Y is an ACTUAL primary blowup, as follows. For L=O(-E), choose finitely many affine nonvanishing section charts of powers of L, equalize their degrees to n, and select finite R-algebra generators of each chart. Stacks28.27.11 section localization writes these generators as t/s\_i\textasciicircum{}k. For one common K, use sections s\_i\textasciicircum{}K and t s\_i\textasciicircum{}(K-k) of L\textasciicircum{}(nK). Their projective chart ratios include every affine generator, hence give an immersion. Properness makes it a closed immersion. All these sections lie in R, since L\textasciicircum{}(nK) is a subsheaf of O. Their ideal J generates L\textasciicircum{}(nK), is proper because E is nonempty, and is unit off the marked point, hence primary. The projective map factors through the integral Rees blowup Bl\_J(R). Its factor is a closed immersion containing the dense punctured open, so its defining ideal is zero. Thus Y=Bl\_J(R). This proves the finite-list refinement actually used, not arbitrary-model cofinality.

Write the exceptional curves E\_i, the pulled-back gauge divisor E\_a=sum a\_i E\_i with a\_i\textgreater{}0, and signed obstacle D=sum d\_i E\_i. Regular local factoriality makes the E\_i actual effective Cartier divisors. Put Q\_ij=deg\_(E\_i) O(E\_j), with residue-field degrees included. Off-diagonal entries are nonnegative and symmetric: the effective zero divisors are computed by the same local intersection lengths. A nonzero section of a line bundle on any proper integral curve has nonnegative degree by its effective zero-divisor sequence. Every contracted integral curve is an E\_i. Thus curve-nefness is Qz\textgreater{}=0.

For c\_i=min(d\_i,0), the feasible vectors z\textless{}=c, Qz\textgreater{}=0 contain the barrier b=-M a for sufficiently large M. The barrier is nef because the negative gauge divisor is globally generated. Coordinatewise maxima preserve the inequalities: for coordinate i keep an input's diagonal term and increase the nonnegative off-diagonal terms. In the compact box b\textless{}=z\textless{}=c, maximize each coordinate and take the finite coordinatewise maximum of those maximizers. This gives z* dominating the box. Every feasible z has max(z,b) in the box, so z* dominates ALL feasible vectors. Z*=sum z\_i*E\_i is ACTUAL negative effective real Cartier on ACTUAL Y, with degree \textgreater{}=0 on EVERY contracted integral curve. Its evaluation is in the geometric core and below f\_D on EVERY support.

For ANY rational primary competitor q=-v(I)/s\textless{}=f\_D let F\_I=sum ord\_(E\_i)(I) E\_i. Normal height-one intersection shows every ideal generator is a global section of O(-F\_I); at the generic point of each E\_i one generator attains the minimum and induces a nonzero section on that CURVE. Consequently A=-F\_I/s is curve-nef, nonpositive, and divisorial tests ord\_(E\_i)/a\_i give A\_i\textless{}=d\_i. Thus A\textless{}=Z*. On EVERY support p, use the valuation ring of Frac(R/p), not a homomorphism from Frac(R). Its generic centre is outside the exceptional fibre, properness gives a centre on Y, and the invertible exceptional lattice has a nonzero generic trivialization. Dividing each ideal generator by its local lattice generator is regular, proving q\textless{}=f\_A\textless{}=f\_(Z*); killed generators have infinite value and cause no exception. The raw rational-competitor identity then proves p(f\_D)=f\_(Z*).

Signed primary model obstacles are uniformly dense: their rational difference lattice obeys min(vI-vJ,vK-vH)=v(IH+KJ)-v(JH), contains constants through a, and separates all semivaluations using (g)+a\textasciicircum{}N, including different supports. Compact lattice Stone-Weierstrass and the raw Lipschitz bound give CE for EVERY continuous signed obstacle on EVERY such normal surface, arbitrary residue field and equal or mixed characteristic. Every canonical degree converges uniformly. Rational coefficient linear programming does not prove ideal attainment, semiampleness, finite generation, or a general O(1/m) rate.

For the characteristic-zero smooth theorem used later, the published BFJ setting is ONLY a smooth affine CLOSED germ over an algebraically closed characteristic-zero field, dimension\textgreater{}=2; DVR and empty cases are direct. For a signed model F=F composed with a fixed retraction, adjoin a negative constant barrier to the max-directed family of nef coefficients below F. Weak closedness gives finite nef W. BFJ3.11 gives its substantive CLOSED-simplex continuous function H, and its ALL-retraction infimum. Divisorial density and refinements of the fixed model show H\textless{}=F on EVERY support. G331's exact BFJ D comparison shows every D competitor is below H. The canonical ideals are L-infinity(mW): to extend their divisorial tests to every support, truncate (g)+mhat\textasciicircum{}N with N\textgreater{}=m max(0,-min F), use closed-simplex continuity and cofinal retractions, and then v(g)\textgreater{}=v((g)+mhat\textasciicircum{}N). G331 contracts these primary ideals exactly to the local ring.

If W=0, H=0. Otherwise BFJ3.14 makes H strictly negative and thinness is bounded on the FIXED compact complex, so A\_thin\textless{}=-CW there. The complete3.18 proof gives W\_t\textless{}=W\textless{}=(1-C/t)W\_(t-C) for real t\textgreater{}C: part(i) applies to any positive real scalar multiple, and the defining thinness argument of(ii) applies to real t. Taking t=m+C gives (1+C/m)W\textless{}=W\_m\textless{}=W. Closed simplices and ALL-retraction infima yield 0\textless{}=H-h\_m\textless{}=C max(0,-min F)/m on EVERY support. Uniform primary approximation proves H in C. Uniform model density and Lipschitz then prove smooth CE for all finite continuous signed obstacles at exactly the published hypotheses. Hironaka smooth char-zero principalization, multiplier subadditivity, thinness and smooth Izumi remain disclosed deeper foundations, not positive/mixed-characteristic premises. The separately retained coefficient-field, quotient and root proofs are DATA; no enlargement of this BFJ application is needed here.

On a reduced excellent curve the normalization centres are DVRs, the regular conductor has finite value at every branch, and every boundary fibre is singleton. Actual real Cartier slopes have empty curve tests; semilocal gluing and G331 positive descent give all nonpositive functions. Hence the geometric envelope is min(f,0), equal to the ideal envelope. This does not identify the boundary cone with literal multibranch classical/CLD radial extensions; the mixed-zero obstruction remains. Nilpotents reduce away and dimension-zero reduction has empty boundary.

Evidence and limitations: all four current416 actual FINALs were physically read with ALL five fields, and all four distinct expanded proofs were separately read in bounded untruncated outputs. Whole current unique preamble/abstract/ledger/verdict/residual and scope records were read. The six immediate surface primary text units and eight SH/BFJ units were physically read whole; both SH formula pixels were viewed. The initial binding inventory and one inherited413 display were truncated and are excluded as whole-reading credit. The missing inherited source range950--986 was then read separately; ranges1039--1107 and1108--1216 completed the retained413 actual/expanded proof units. Exact literal repetitions carry identity, not extra independent-proof credit. The four supplied native programs were physically read as inert DATA and never executed/imported. No complete fresh recursive historical/whole-book audit is claimed. The source/companion are retained entire. Artifacts n1.md and n1.tex contain this reconstruction; rerun own programs: python assemble.py; python verify.py. Execution is artifact consistency only.

\subsection*{Expanded canonical and surface formulae}
\[
A_m(f)=\{g:v(g)+mf(v)\geq0\ \forall v\},\quad
h_m=-m^{-1}v(A_m),\quad p(f)=\sup_m h_m,
\]
\[
H_N-\frac{NM_1}{m}\leq h_m\leq p(f),\qquad
p(f)\text{ continuous}\Longleftrightarrow h_m\to p(f)\text{ uniformly}.
\]
On a proved regular actual primary surface model,
\[
z^*=\max\{z:z\leq\min(d,0),\ Qz\geq0\},\qquad
p(f_D)=f_{\sum_i z_i^*E_i}.
\]
The maximum denotes the greatest vector, proved by compact coordinate
maximization and closure under coordinatewise maxima, not an assumed
linear-programming output. The every-support ideal comparison is
$-v(I)/s\leq f_{-F_I/s}\leq f_{Z^*}$.

